% =============================================================================
% Section 6: Discussion and Lessons
%
% Intent: three threads.
%   (1) Why agentic pipelines underperformed RAG in our setup. Hallucination
%       amplification, tool-call overhead, evaluation against a static
%       snapshot. Concrete: the "always provide your best guess" prompt
%       pattern as a known anti-pattern.
%   (2) What the paradoxical-completeness pattern (high completeness,
%       low correctness) reveals about LLM-as-judge for domain ops.
%       Forward reference back to Sec. 5.7.
%   (3) Deployment lessons: latency vs. throughput tradeoffs, cost,
%       prompt fragility. The MoE-vs-dense throughput observation
%       (gemma4 vs qwen3) belongs here.
%
% Budget: 0.5 pp. Tight. No padding.
% =============================================================================

\section{Discussion and Lessons}
\label{sec:discussion}

% TODO: 3 paragraphs.
%
% P1 --- agentic regression. Why the Copilot agent scored lowest on
%        correctness despite having tools. Three causes: hallucination
%        amplification, tool-call overhead consuming context, evaluation
%        runs against a static snapshot of doc+tool state. Concrete remedy:
%        revise the agent prompt to permit clarification questions and to
%        express uncertainty.
%
% P2 --- judge fragility. The paradoxical-completeness pattern means the
%        judge rewards confident structure independently of factual
%        accuracy. The methodology's response: reference-free with a
%        production-answer ceiling, multiple judges with agreement metrics.
%        Open question for the field: judge calibration on domain ops.
%
% P3 --- deployment lessons. Latency-vs-correctness tradeoffs we observed.
%        Cost of serving 32B vs 120B locally. Prompt fragility under tool
%        addition. The MoE throughput finding (gemma4 ~3x faster than
%        qwen3:32b on identical hardware) and what it implies for
%        latency-sensitive ops deployments.
